\chapter{$U_\theta$-模的局部化}
\label{ch:03.3}

\section{$U_\theta$-模的局部化}
\label{sec:03.3.1}

设 $\lambda\in\h^*$，$\theta$ 是相应的 Weyl 群轨道。于是可以定义右正合\index{youzhenghe@右正合|hyperpage}共变函子 $\Delta_\lambda:\mathcal{M}(U_\theta)\to\mathcal{M}_{qc}(\D_\lambda)$ 为
$$\Delta_\lambda(V)=\D_\lambda\otimes_{U_\theta}V,\qquad V\in\mathcal{M}(U_\theta).$$
称之为\emph{局部化函子\index{jubuhuahanzi@局部化函子|hyperpage}}（localization functor）。由于对任意 $W\in\mathcal{M}(\D_\lambda)$ 有
$$\Gamma(X,W)=\Hom_{\D_\lambda}(\D_\lambda,W),$$
可知 $\Delta_\lambda$ 是整体截面函子 $\Gamma$ 的左伴随\index{zuobansui@左伴随|hyperpage}函子\index{bansuihanzi@伴随函子|hyperpage}，即
$$\Hom_{\D_\lambda}(\Delta_\lambda(V),W)=\Hom_{U_\theta}(V,\Gamma(X,W)),$$
其中 $V\in\mathcal{M}(U_\theta)$，$W\in\mathcal{M}(\D_\lambda)$。

特别地，存在从恒等函子到 $\Gamma\circ\Delta_\lambda$ 的函子态射 $\phi$。对任意 $V\in\mathcal{M}(U_\theta)$，它由自然态射
$$\phi_V:V\longrightarrow\Gamma(X,\Delta_\lambda(V))$$
给出。

先设 $\lambda\in\h^*$ 反支配。

\begin{lemma}\label{lem:03.1.1}
设 $\lambda\in\h^*$ 反支配。则 $V$ 到 $\Gamma(X,\Delta_\lambda(V))$ 的自然映射 $\phi_V$ 是 $\g$-模同构。
\end{lemma}

\begin{proof}
若 $V=U_\theta$，这由 C.6.1 得出。又由 C.3.2 知，此时 $\Gamma$ 正合。这意味着 $\Gamma\circ\Delta_\lambda$ 右正合。设
$$(U_\theta)^{(J)}\longrightarrow(U_\theta)^{(I)}\longrightarrow V\longrightarrow 0$$
是 $\g$-模的正合列。则有交换图表
\begin{center}
{\small
\begin{tikzcd}[column sep=small]
(U_\theta)^{(J)} \arrow[r] \arrow[d] & (U_\theta)^{(I)} \arrow[r] \arrow[d] & V \arrow[r] \arrow[d] & 0 \\
\Gamma(X,\Delta_\lambda(U_\theta))^{(J)} \arrow[r] & \Gamma(X,\Delta_\lambda(U_\theta))^{(I)} \arrow[r] & \Gamma(X,\Delta_\lambda(V)) \arrow[r] & 0
\end{tikzcd}}
\end{center}
其两行正合，且前两个竖直箭头是同构。由此可知第三个竖直箭头也是同构。
\end{proof}

\begin{annote}
底本中形如 C.6.1、C.3.2 的编号，指向参考文献 [9]、[10] 所载的两部未刊讲义（\emph{Lectures on algebraic D-modules} 与 \emph{Lectures on derived categories}）；本书没有相应的附录，故一律照录原编号，不改写成内部引用。[来源：底本 p.35，参考文献 [9]、[10]]
\end{annote}

另一方面，伴随性还给出从 $\Delta_\lambda\circ\Gamma$ 到恒等函子的函子态射 $\psi$。对任意 $V\in\mathcal{M}(\D_\lambda)$，它由自然态射
$$\psi_V:\Delta_\lambda(\Gamma(X,V))=\D_\lambda\otimes_{U_\theta}\Gamma(X,V)\longrightarrow V$$
给出。设 $V\in\mathcal{M}_{qc}(\D_\lambda)$，$K$ 与 $C$ 分别是 $\psi_V$ 的核与余核\index{yuhe@余核|hyperpage}。则有拟凝聚 $\D_\lambda$-模的正合列
$$0\longrightarrow K\longrightarrow\Delta_\lambda(\Gamma(X,V))\longrightarrow V\longrightarrow C\longrightarrow 0.$$
对它作用 $\Gamma$，并用 C.3.2，得到正合列
$$0\longrightarrow\Gamma(X,K)\longrightarrow\Gamma(X,\Delta_\lambda(\Gamma(X,V)))\longrightarrow\Gamma(X,V)\longrightarrow\Gamma(X,C)\longrightarrow 0.$$
于是由 \ref{lem:03.1.1} 可见 $\Gamma(X,K)=0$，$\Gamma(X,C)=0$。由此得到下面的结果。

记 $Q\mathcal{M}_{qc}(\D_\lambda)$ 为 $\mathcal{M}_{qc}(\D_\lambda)$ 关于「整体截面为零的拟凝聚 $\D_\lambda$-模」所成子范畴的商范畴\index{shangfanchou@商范畴|hyperpage} [6, 1.11]。设 $Q$ 是 $\mathcal{M}_{qc}(\D_\lambda)$ 到 $Q\mathcal{M}_{qc}(\D_\lambda)$ 的商函子。显然，$\Gamma$ 诱导出从 $Q\mathcal{M}_{qc}(\D_\lambda)$ 到 $\mathcal{M}(U_\theta)$ 的正合函子；以下滥用记号，仍记作 $\Gamma$。

\begin{theorem}\label{thm:03.1.2}
设 $\lambda\in\h^*$ 反支配。则函子 $Q\circ\Delta_\lambda:\mathcal{M}(U_\theta)\to Q\mathcal{M}_{qc}(\D_\lambda)$ 是范畴等价，其逆为 $\Gamma$。
\end{theorem}

若 $\lambda$ 反支配且正则，则由 C.4.4，$\mathcal{M}_{qc}(\D_\lambda)$ 中每个对象都可由其整体截面生成。因此此时 $Q\mathcal{M}_{qc}(\D_\lambda)=\mathcal{M}_{qc}(\D_\lambda)$。

\begin{corollary}\label{cor:03.1.3}
设 $\lambda\in\h^*$ 反支配且正则。则函子 $\Delta_\lambda:\mathcal{M}(U_\theta)\to\mathcal{M}_{qc}(\D_\lambda)$ 是范畴等价，其逆为 $\Gamma$。
\end{corollary}

作为这个范畴等价的第一个应用，我们来证明关于环 $U_\theta$ 的同调维数的一个结果。

\begin{theorem}\label{thm:03.1.4}
设 $\theta$ 是 $\h^*$ 中由正则元素组成的 Weyl 群轨道。则 $U_\theta$ 的同调维数满足
$$\operatorname{hd}(U_\theta)\le\dim X+\frac{1}{2}\operatorname{Card}\Sigma_\lambda.$$
\end{theorem}

\begin{proof}
设 $\lambda\in\theta$ 反支配。由 [1, VI.1.10(ii)] 知，$\D_{X,x}$ 的同调维数等于 $\dim X$。由于 $\D_\lambda$ 是扭微分算子层，可知 $\operatorname{hd}(\D_{\lambda,x})=\dim X$。此外，由 [6, 4.2.2]，对任意 $i\in\Z_+$、$V\in\mathcal{M}_{coh}(\D_\lambda)$ 与 $U\in\mathcal{M}_{qc}(\D_\lambda)$ 有
$$\Ext^i_{\D_\lambda}(V,U)_x=\Ext^i_{\D_{\lambda,x}}(V_x,U_x).$$
于是当 $i>\dim X$ 时 $\Ext^i_{\D_\lambda}(V,U)=0$。

另一方面，有谱序列
$$H^p(X,\Ext^q_{\D_\lambda}(V,U))\Longrightarrow\Ext^{p+q}_{\D_\lambda}(V,U)$$
（[6, 4.2.1]）；而任意阿贝尔群层的上同调在次数大于 $\dim X$ 时为零（[8, III.2.7]）。由此可见，当 $i>2\dim X$ 时 $\Ext^i_{\D_\lambda}(V,U)=0$。现在由 \ref{cor:03.1.3}，$\mathcal{M}(U_\theta)$ 与 $\mathcal{M}_{qc}(\D_\lambda)$ 等价，且任意有限生成 $U_\theta$-模的局部化都属于 $\mathcal{M}_{coh}(\D_\lambda)$。这意味着，对任意 $V\in\mathcal{M}_{fg}(U_\theta)$ 与 $U\in\mathcal{M}(U_\theta)$，当 $i>2\dim X$ 时 $\Ext^i_{U_\theta}(V,U)=0$。由 [3, Ch.~X, \S8, no.~3, Cor.~of Prop.~4] 可知 $\operatorname{hd}(U_\theta)\le 2\dim X$。

若 $\mathcal{V}$ 是任意凝聚 $\D_\lambda$-模，则 $\Ext^i_{\D_\lambda}(\mathcal{V},\D_\lambda)$（$i\in\Z_+$）是凝聚 $\D_{-\lambda}$-模 [9]。由于 $\lambda$ 正则反支配，$n(-\lambda)=\frac{1}{2}\operatorname{Card}\Sigma_\lambda\le\dim X$；又由 C.3.1，$\Ext^i_{\D_\lambda}(\mathcal{V},\D_\lambda)$ 的上同调在 $\frac{1}{2}\operatorname{Card}\Sigma_\lambda$ 以上消没。因此，与前面的论证一样，可得当 $i>\dim X+\frac{1}{2}\operatorname{Card}\Sigma_\lambda$ 时 $\Ext^i_{\D_\lambda}(\mathcal{V},\D_\lambda)=0$。由范畴等价立即可见，对任意 $V\in\mathcal{M}_{fg}(U_\theta)$ 与 $i>\dim X+\frac{1}{2}\operatorname{Card}\Sigma_\lambda$ 有 $\Ext^i_{U_\theta}(V,U_\theta)=0$。

现在设 $V,U\in\mathcal{M}_{fg}(U_\theta)$。由于 $U_\theta$ 是左诺特的，有正合列
$$0\longrightarrow U'\longrightarrow U_\theta^p\longrightarrow U\longrightarrow 0,$$
其中 $U'$ 是有限生成 $U_\theta$-模。由 $\Ext^\bullet_{U_\theta}(V,-)$ 相应的长正合列可见，当 $j>\dim X+\frac{1}{2}\operatorname{Card}\Sigma_\lambda$ 时，连接态射
$$\Ext^j_{U_\theta}(V,U)\longrightarrow\Ext^{j+1}_{U_\theta}(V,U')$$
是同构。由于 $U_\theta$ 的同调维数有限，对 $j$ 作倒向归纳可知，当 $j>\dim X+\frac{1}{2}\operatorname{Card}\Sigma_\lambda$ 时 $\Ext^j_{U_\theta}(V,U)=0$。再由 [3, Ch.~X, \S8, no.~3, Cor.~of Prop.~4] 即得
$$\operatorname{hd}(U_\theta)\le\dim X+\frac{1}{2}\operatorname{Card}\Sigma_\lambda.$$
\end{proof}

\begin{remark}\label{rem:03.1.5}
以后在 2.8 中会看到：与 \ref{thm:03.1.4} 相反，若 $\theta$ 含有 $\h^*$ 的奇异元素，则 $U_\theta$ 的同调维数无穷。

此外，对正则整权 $\lambda$ 的任一 $W$-轨道 $\theta$，上述关于 $U_\theta$ 同调维数的估计是精确的。为看出这一点，设 $\lambda\in P(\Sigma)$ 正则反支配，$F$ 是以 $\lambda+\rho$ 为最低权的不可约有限维 $\g$-模。则由 Borel--Weil--Bott 定理与 \ref{thm:03.1.2} 知 $\Delta_\lambda(F)=\calO(\lambda+\rho)$。于是由 ……，有
$$\Ext^p_{\D_\lambda}(\calO(\lambda+\rho),\D_\lambda)=0,\qquad p\ne\dim X,$$
且作为左 $\D_{-\lambda}$-模有
$$\Ext^{\dim X}_{\D_\lambda}(\calO(\lambda+\rho),\D_\lambda)=\calO(-\lambda+\rho).$$
因此，再次应用 Borel--Weil--Bott 定理，可见
$$H^p(X,\Ext^{\dim X}_{\D_\lambda}(\calO(\lambda+\rho),\D_\lambda))=0,\qquad p\ne\dim X,$$
且
$$H^{\dim X}(X,\Ext^{\dim X}_{\D_\lambda}(\calO(\lambda+\rho),\D_\lambda))\ne 0.$$
由联系 $\Ext_{U_\theta}$ 与 $\Ext_{\D_\lambda}$ 的 Grothendieck 谱序列，这意味着当 $p\ne 2\dim X$ 时
$$\Ext^p_{U_\theta}(F,U_\theta)=\Ext^p_{\D_\lambda}(\calO(\lambda+\rho),\D_\lambda)=0,$$
且
$$\Ext^{2\dim X}_{U_\theta}(F,U_\theta)=\Ext^{2\dim X}_{\D_\lambda}(\calO(\lambda+\rho),\D_\lambda)\ne 0.$$
因此此时 $\operatorname{hd}(U_\theta)=2\dim X=\operatorname{Card}\Sigma$。
\end{remark}

\begin{annote}
这一段底本有两处引文问题。其一，「于是由 ……」处原文引文空缺（印作「by ...」）；同类空缺还见于 \ref{lem:03.1.19} 的证明，共三处，均照录为「由 ……」。其二，「联系 $\Ext_{U_\theta}$ 与 $\Ext_{\D_\lambda}$ 的 Grothendieck 谱序列」底本印作两端都是 $\Ext_{\D_\lambda}$，与紧随其后的等式不合，已按上下文补正。[来源：底本 p.37 与 p.41，上下文自证]
\end{annote}

作为第二个应用，我们考虑 $X$ 上 $\D_\lambda$-模的各种导出范畴。仍如前面，记 $D(\mathcal{M}(\D_\lambda))$ 为 $\D_\lambda$-模的导出范畴，$D_{qc}(\mathcal{M}(\D_\lambda))$ 为其中由上同调拟凝聚的复形组成的满子范畴。也可以考虑拟凝聚 $\D_\lambda$-模的导出范畴 $D(\D_\lambda)=D(\mathcal{M}_{qc}(\D_\lambda))$。如前所述，对任意 $\mu\in P(\Sigma)$，几何平移函子 $V\mapsto V(\mu)$ 是 $\mathcal{M}(\D_\lambda)$ 与 $\mathcal{M}(\D_{\lambda+\mu})$ 之间的等价，它也诱导子范畴 $\mathcal{M}_{qc}(\D_\lambda)$ 与 $\mathcal{M}_{qc}(\D_{\lambda+\mu})$ 的等价。此外，它还诱导相应导出范畴 $D(\mathcal{M}(\D_\lambda))$、$D_{qc}(\mathcal{M}(\D_\lambda))$、$D(\D_\lambda)$ 分别与 $D(\mathcal{M}(\D_{\lambda+\mu}))$、$D_{qc}(\mathcal{M}(\D_{\lambda+\mu}))$、$D(\D_{\lambda+\mu})$ 的等价。又，从 $D(\D_\lambda)$ 到 $D_{qc}(\mathcal{M}(\D_\lambda))$ 的典范函子 $\Phi_\lambda$ 关于这些平移函子满足自然的交换性。

\begin{theorem}\label{thm:03.1.6}
函子 $\Phi_\lambda:D(\D_\lambda)\to D_{qc}(\mathcal{M}(\D_\lambda))$ 是范畴等价。
\end{theorem}

\begin{proof}
首先，由前述讨论，经平移可设 $\lambda$ 反支配且正则。此时局部化函子 $\Delta_\lambda$ 正合。因此它诱导出从 $U_\theta$-模的导出范畴 $D(U_\theta)$ 到 $D(\D_\lambda)$ 的函子
$$\Delta_\lambda:V^\bullet\longmapsto\D_\lambda\otimes_{U_\theta}V^\bullet.$$
另一方面，由 C.3.1，$\mathcal{M}_{qc}(\D_\lambda)$ 中每个对象都是 $\Gamma$-零调的；结合 \ref{cor:03.1.3} 立即推出，函子 $\Gamma:D(\D_\lambda)\to D(U_\theta)$ 是范畴等价，其逆为局部化函子 $\Delta_\lambda$。此外，由 $\Gamma$ 在 $\mathcal{M}(\D_\lambda)$ 上右上同调维数有限，有导出函子 $R\Gamma:D_{qc}(\mathcal{M}(\D_\lambda))\to D(U_\theta)$，且显然 $R\Gamma\circ\Phi_\lambda=\Gamma$。又，任一 $C^\bullet\in D_{qc}(\mathcal{M}(\D_\lambda))$ 都可换成由 $\mathcal{M}(\D_\lambda)$ 中 $\Gamma$-零调对象组成的拟同构复形 $A^\bullet$。于是 $R\Gamma(C^\bullet)=\Gamma(A^\bullet)$。置 $D^\bullet=\D_\lambda\otimes_{U_\theta}\Gamma(A^\bullet)$，有自然同态 $\varphi:D^\bullet\to A^\bullet$。我们断言它是拟同构。显然，由 $\varphi$ 的定义与 \ref{lem:03.1.1}，$R\Gamma(\varphi)$ 是拟同构。因此由 C.4.1 可见 $\varphi$ 是拟同构。于是 $\Phi_\lambda\circ(\Delta_\lambda\circ R\Gamma)$ 同构于 $D_{qc}(\mathcal{M}(\D_\lambda))$ 上的恒等函子。另一方面，$(\Delta_\lambda\circ R\Gamma)\circ\Phi_\lambda$ 同构于 $D(\D_\lambda)$ 上的恒等函子。因此 $\Phi_\lambda$ 是范畴等价。
\end{proof}

对上有界与下有界复形的导出范畴，类似的陈述同样成立。

设 $\lambda\in\h^*$，$\theta=W\cdot\lambda$。记 $D^-(U_\theta)$ 为 $U_\theta$-模的上方有界复形所成的导出范畴。定义局部化函子 $L\Delta_\lambda:D^-(U_\theta)\to D(\D_\lambda)$ 为
$$L\Delta_\lambda(V^\bullet)=\D_\lambda\overset{L}{\otimes}_{U_\theta}V^\bullet,\qquad V^\bullet\in D^-(U_\theta).$$
若 $\lambda$ 正则，则由 \ref{thm:03.1.4} 知，局部化函子 $\Delta_\lambda$ 的左上同调维数不超过 $2\dim X$。因此可以把 $L\Delta_\lambda$ 延拓为从 $D(U_\theta)$ 到 $D(\D_\lambda)$ 的函子。

\begin{lemma}\label{lem:03.1.7}
设 $P\in\mathcal{M}(U_\theta)$ 是投射模\index{toushemo@投射模|hyperpage}。则其局部化 $\Delta_\lambda(P)$ 是 $\Gamma$-零调的，且态射 $\phi_P:P\to\Gamma(X,\Delta_\lambda(P))$ 是同构。
\end{lemma}

\begin{proof}
由 C.6.1 知，该结论对自由 $U_\theta$-模成立；而任意投射 $U_\theta$-模都是自由 $U_\theta$-模的直和项。
\end{proof}

设 $V^\bullet\in D^-(U_\theta)$。则存在由投射 $U_\theta$-模组成的复形 $P^\bullet\in D(U_\theta)$ 与拟同构 $\alpha_{V^\bullet}:P^\bullet\to V^\bullet$，且 $L\Delta_\lambda(V^\bullet)=\Delta_\lambda(P^\bullet)$。由 \ref{lem:03.1.7} 可见，存在从 $P^\bullet$ 到 $\Gamma(\Delta_\lambda(P^\bullet))=R\Gamma(\Delta_\lambda(P^\bullet))$ 的自然同构。由此得到下面的结果。

\begin{lemma}\label{lem:03.1.8}
函子 $R\Gamma\circ L\Delta_\lambda$ 从 $D^-(U_\theta)$ 到自身，且同构于 $D^-(U_\theta)$ 上的恒等函子。
\end{lemma}

设 $D:\mathcal{M}(U_\theta)\to D^-(U_\theta)$ 是把 $V\in\mathcal{M}(U_\theta)$ 映为复形 $D(V)\in D^-(U_\theta)$ 的函子，该复形除 $0$ 次处等于 $V$ 外，其余各次均为零。若对任意 $V\in\mathcal{M}(U_\theta)$，记 $L^j\Delta_\lambda(V)$ 为复形 $L\Delta_\lambda(D(V))$ 的第 $j$ 个上同调，则得到函子 $V\mapsto L^j\Delta_\lambda(V)$，从 $\mathcal{M}(U_\theta)$ 到 $\mathcal{M}_{qc}(\D_\lambda)$，它正是 $\Delta_\lambda$ 的第 $(-j)$ 个左导出函子。因此由 \ref{lem:03.1.8} 得到下面的结果。

\begin{corollary}\label{cor:03.1.9}
设 $V\in\mathcal{M}(U_\theta)$。则存在以
$$E_2^{p,q}=H^p(X,L^q\Delta_\lambda(V))$$
为 $E_2$ 项、收敛到 $V$ 的上同调谱序列。
\end{corollary}

\begin{corollary}\label{cor:03.1.10}
设 $F\in\mathcal{M}(U_\theta)$ 是平坦\index{pingtan@平坦|hyperpage}模。则其局部化 $\Delta_\lambda(F)$ 是 $\Gamma$-零调的，且态射 $\phi_F:F\to\Gamma(X,\Delta_\lambda(F))$ 是同构。
\end{corollary}

\begin{proof}
按定义，当 $q\ne 0$ 时 $L^q\Delta_\lambda(F)=0$。因此 \ref{cor:03.1.9} 中的谱序列退化，于是当 $i>0$ 时 $H^i(X,\Delta_\lambda(F))=0$，且 $\phi_F:F\to\Gamma(X,\Delta_\lambda(F))$ 是同构。
\end{proof}

现在设 $\lambda\in\h^*$ 正则。则 $U_\theta$ 的同调维数有限，因此任一 $V^\bullet\in D(U_\theta)$ 都拟同构于由投射 $U_\theta$-模组成的复形 $P^\bullet\in D(U_\theta)$ [10]。再次利用 \ref{lem:03.1.7}，可以证明 \ref{lem:03.1.8} 的下述变体。

\begin{lemma}\label{lem:03.1.11}
设 $\lambda\in\h^*$ 正则。则函子 $R\Gamma\circ L\Delta_\lambda$ 从 $D(U_\theta)$ 到自身，且同构于 $D(U_\theta)$ 上的恒等函子。
\end{lemma}

最后得到 \ref{cor:03.1.3} 的下述类似结果。

\begin{theorem}\label{thm:03.1.12}
设 $\lambda\in\h^*$ 正则。则函子 $L\Delta_\lambda:D(U_\theta)\to D(\D_\lambda)$ 是范畴等价，其逆为 $R\Gamma$。
\end{theorem}

\begin{proof}
设 $V^\bullet\in D_{qc}(\mathcal{M}(\D_\lambda))$。则存在由 $\Gamma$-零调 $\D_\lambda$-模组成的复形 $C^\bullet\in D_{qc}(\mathcal{M}(\D_\lambda))$ 与拟同构 $\beta_{V^\bullet}:V^\bullet\to C^\bullet$，且 $R\Gamma(V^\bullet)=\Gamma(C^\bullet)$。又存在由投射 $U_\theta$-模组成的复形 $P^\bullet\in D(U_\theta)$ 与拟同构 $\alpha_{\Gamma(C^\bullet)}:P^\bullet\to\Gamma(C^\bullet)$，且 $L\Delta_\lambda(\Gamma(C^\bullet))=\Delta_\lambda(P^\bullet)$。于是得到从 $L\Delta_\lambda(R\Gamma(V^\bullet))=\Delta_\lambda(P^\bullet)$ 到 $C^\bullet$ 的自然态射。由此得到从 $L\Delta_\lambda\circ R\Gamma$ 到 $D_{qc}(\mathcal{M}(\D_\lambda))$ 上恒等函子的函子态射。由 \ref{thm:03.1.6}，与 $\Phi_\lambda$ 复合后得到从函子 $L\Delta_\lambda\circ R\Gamma$ 到 $D(\D_\lambda)$ 上恒等函子的态射 $\psi$。于是对任意复形 $V^\bullet\in D(\D_\lambda)$，存在态射 $\psi_{V^\bullet}:L\Delta_\lambda(R\Gamma(V^\bullet))\to V^\bullet$；检查其定义并用 \ref{lem:03.1.11}，可见 $R\Gamma(\psi_{V^\bullet})$ 是拟同构。于是由 C.4.1 可知，$\psi$ 是函子同构。
\end{proof}

特别地，这意味着 $L\Delta_\lambda$ 是 $D^b(U_\theta)$ 与 $D^b(\D_\lambda)$ 之间的范畴等价，$R\Gamma$ 是其逆。

\begin{theorem}\label{thm:03.1.13}
设 $\lambda\in\h^*$ 正则。则 $\Delta_\lambda$ 的左上同调维数不超过 $n(\lambda)$。
\end{theorem}

\begin{proof}
设 $V^\bullet\in D(\D_\lambda)$，$k\in\Z$。则截断复形
$$\sigma_{\le k}(V^\bullet):\quad\cdots\longrightarrow V_p\longrightarrow\cdots\longrightarrow V_{k-2}\longrightarrow V_{k-1}\longrightarrow\ker d_k\longrightarrow 0\longrightarrow\cdots$$
自然映入 $V^\bullet$，且这个复形态射对 $p\le k$ 诱导同构 $H^p(\sigma_{\le k}(V^\bullet))\to H^p(V^\bullet)$。设 $V\in\mathcal{M}(U_\theta)$，$V^\bullet=L\Delta_\lambda(D(V))$，并设 $-k>n(\lambda)$。则有上同调谱序列
$$H^p(X,H^q(\sigma_{\le k}(V^\bullet)))\Longrightarrow H^{p+q}(R\Gamma(\sigma_{\le k}(V^\bullet)))$$
（[6], II.2.4）。由 C.3.1 可知，当 $p\in\Z_+$ 时 $H^p(R\Gamma(\sigma_{\le k}(V^\bullet)))=0$。于是由 \ref{thm:03.1.12}，
\begin{align*}
\Hom_{D(\D_\lambda)}(\sigma_{\le k}(V^\bullet),V^\bullet)
&=\Hom_{D(U_\theta)}(R\Gamma(\sigma_{\le k}(V^\bullet)),R\Gamma(V^\bullet))\\
&=\Hom_{D(U_\theta)}(R\Gamma(\sigma_{\le k}(V^\bullet)),D(V))=0.
\end{align*}
因此当 $p\le -k$ 时 $L^p\Delta_\lambda(V)=H^p(V^\bullet)=0$。
\end{proof}

\begin{remark}\label{rem:03.1.14}
相反，以后在 2.7 中会看到，若 $\lambda$ 奇异，则 $\Delta_\lambda$ 的左上同调维数无穷。
\end{remark}

设 $\mathcal{A}$ 是阿贝尔范畴，$D^b(\mathcal{A})$ 是其有界复形的导出范畴。对 $s\in\Z$，有截断函子 $\tau_{\ge s}$ 与 $\tau_{\le s}$，从 $D^b(\mathcal{A})$ 到自身 [10]。若 $A^\bullet$ 是 $D^b(\mathcal{A})$ 中的复形，则 $\tau_{\ge s}(A^\bullet)$ 是如下复形：次数小于 $s$ 处为零，$\tau_{\ge s}(A^\bullet)^s=\operatorname{coker} d^{s-1}$，且当 $q>s$ 时 $\tau_{\ge s}(A^\bullet)^q=A^q$，微分由 $A^\bullet$ 的微分诱导。另一方面，$\tau_{\le s}(A^\bullet)$ 是如下复形：次数大于 $s$ 处为零，$\tau_{\le s}(A^\bullet)^s=\ker d^s$，且当 $q<s$ 时 $\tau_{\le s}(A^\bullet)^q=A^q$，微分由 $A^\bullet$ 的微分诱导。自然态射 $\tau_{\le s}(A^\bullet)\to A^\bullet$ 与 $A^\bullet\to\tau_{\ge s}(A^\bullet)$ 分别在次数 $\le s$ 与 $\ge s$ 处诱导上同调的同构。

\begin{lemma}\label{lem:03.1.15}
设 $C^\bullet,D^\bullet\in D^b(\mathcal{A})$，且
\begin{enumerate}[label=(\roman*),leftmargin=2.6em,itemsep=0.1em]
\item 当 $q>0$ 时 $H^q(C^\bullet)=0$；
\item 当 $q<0$ 时 $H^q(D^\bullet)=0$。
\end{enumerate}
则
$$H^0:\Hom_{D^b(\mathcal{A})}(C^\bullet,D^\bullet)\longrightarrow\Hom_{\mathcal{A}}(H^0(C^\bullet),H^0(D^\bullet))$$
是同构。
\end{lemma}

\begin{proof}
由假设，$\tau_{\le 0}(C^\bullet)\to C^\bullet$ 与 $D^\bullet\to\tau_{\ge 0}(D^\bullet)$ 都是拟同构；把它们与 $\varphi$ 复合，可设当 $q>0$ 时 $C^q=0$，当 $q<0$ 时 $D^q=0$。于是 $\Hom_{\mathcal{A}}(H^0(C^\bullet),H^0(D^\bullet))$ 的每个元素都定义了复形 $C^\bullet$ 到复形 $D^\bullet$ 的态射，故我们的映射是满的。为证单性，考虑态射 $\varphi\in\Hom_{D^b(\mathcal{A})}(C^\bullet,D^\bullet)$，使 $H^0(\varphi)=0$。按导出范畴中态射的定义，存在复形 $B^\bullet\in D^b(\mathcal{A})$ 与复形态射 $q:B^\bullet\to C^\bullet$、$f:B^\bullet\to D^\bullet$ 代表 $\varphi$，其中 $q$ 是拟同构。把它们与截断态射 $\tau_{\le 0}(B^\bullet)\to B^\bullet$ 复合，可进一步设 $B^\bullet$ 满足：当 $q>0$ 时 $B^q=0$。但这意味着当 $q\ne 0$ 时 $f^q=0$，且 $\operatorname{im} f^0\subset\ker d^0$、$\operatorname{im} d^{-1}\subset\ker f^0$。于是 $H^0(\varphi)=0$ 推出 $f^0=0$。
\end{proof}

下面的结果是 \ref{lem:03.1.1} 对任意正则 $\lambda\in\h^*$ 的弱推广。

\begin{lemma}\label{lem:03.1.16}
设 $\lambda\in\h^*$ 正则，$\theta=W\cdot\lambda$。设 $V$ 是 $U_\theta$-模，
$$p=\min\{q\in\Z\mid L^{-q}\Delta_\lambda(V)\ne 0\}.$$
若对 $q<p$ 有 $H^q(X,L^{-p}\Delta_\lambda(V))=0$，则存在从 $V$ 到 $H^p(X,L^{-p}\Delta_\lambda(V))$ 的非零态射。
\end{lemma}

\begin{proof}
考虑截断态射
$$L\Delta_\lambda(D(V))\longrightarrow\tau_{\ge -p}\bigl(L\Delta_\lambda(D(V))\bigr)=D\bigl(L^{-p}\Delta_\lambda(V)\bigr)[p].$$
由导出范畴的等价，它给出从 $D(V)$ 到 $R\Gamma\bigl(D(L^{-p}\Delta_\lambda(V))[p]\bigr)$ 的非零态射 $\varphi$。除 $0$ 次外，它在两个复形的上同调模之间诱导零态射；在 $0$ 次处则得到从 $V$ 到 $H^p(X,L^{-p}\Delta_\lambda(V))$ 的态射。由于 $L^{-p}\Delta_\lambda(V)$ 的上同调模在次数 $p$ 以下消没，复形 $R\Gamma\bigl(D(L^{-p}\Delta_\lambda(V))[p]\bigr)$ 满足条件 (ii)。于是由 \ref{lem:03.1.15}，从 $V$ 到
$$R\Gamma\bigl(D(L^{-p}\Delta_\lambda(V))\bigr)[p]^0=H^p(X,L^{-p}\Delta_\lambda(V))$$
的态射 $H^0(\varphi)$ 非零。
\end{proof}

下面讨论一些有限性结果。设 $\mathcal{M}_{fg}(U_\theta)$ 是 $\mathcal{M}(U_\theta)$ 中由有限生成 $U_\theta$-模组成的满子范畴。显然，对任意 $\lambda\in\theta$，$V\in\mathcal{M}_{fg}(U_\theta)$ 的局部化是凝聚 $\D_\lambda$-模。反之，有下面的结果。

\begin{lemma}\label{lem:03.1.17}
设 $\lambda\in\h^*$ 反支配且正则。则对任意 $\mathcal{V}\in\mathcal{M}_{coh}(\D_\lambda)$，$U_\theta$-模 $\Gamma(X,\mathcal{V})$ 是有限生成的。
\end{lemma}

\begin{proof}
设 $\mathcal{V}$ 是凝聚 $\D_\lambda$-模，$V=\Gamma(X,\mathcal{V})$。设 $(V_n;n\in\N)$ 是 $V$ 的一列递增的有限生成 $U_\theta$-子模。对它作局部化，得到 $\mathcal{V}$ 的一列递增的凝聚 $\D_\lambda$-子模 $(\mathcal{V}_n=\Delta_\lambda(V_n);n\in\N)$。由于 $\D_\lambda$ 是诺特环层，序列 $(\mathcal{V}_n;n\in\N)$ 稳定。作用 $\Gamma$ 并用 \ref{cor:03.1.3}，原序列 $(V_n;n\in\N)$ 也稳定。这意味着 $V$ 有限生成。
\end{proof}

由此得到 \ref{cor:03.1.3} 的下述精细化。

\begin{proposition}\label{prop:03.1.18}
设 $\lambda\in\h^*$ 反支配且正则。则函子 $\Delta_\lambda:\mathcal{M}_{fg}(U_\theta)\to\mathcal{M}_{coh}(\D_\lambda)$ 是范畴等价，其逆为 $\Gamma$。
\end{proposition}

现在要把这些结果推广到正则的 $\lambda\in\theta$。首先，$\mathcal{M}_{fg}(U_\theta)$ 是 $\mathcal{M}(U_\theta)$ 的厚子范畴，因此可以考虑 $D^b(U_\theta)$ 的满子范畴 $D^b_{fg}(U_\theta)$，它由所有上同调模有限生成的有界复形组成。

\begin{lemma}\label{lem:03.1.19}
自然函子 $i:D^b(\mathcal{M}_{fg}(U_\theta))\to D^b_{fg}(U_\theta)$ 是范畴等价。
\end{lemma}

\begin{proof}
先证
$$i:D^b(\mathcal{M}_{fg}(U_\theta))\longrightarrow D^b_{fg}(U_\theta)$$
是全忠实的。设 $A^\bullet,B^\bullet\in D^b(\mathcal{M}_{fg}(U_\theta))$，$\phi\in\Hom_{D^b(U_\theta)}(A^\bullet,B^\bullet)$。则存在复形 $C^\bullet\in D^b_{fg}(U_\theta)$、拟同构 $s:C^\bullet\to A^\bullet$ 以及复形态射 $f:C^\bullet\to B^\bullet$，它们定义了 $\phi$。由 …… 可找到复形 $D^\bullet\in D^b(\mathcal{M}_{fg}(U_\theta))$ 与拟同构 $s':D^\bullet\to C^\bullet$。于是 $D^\bullet$、$s\circ s'$ 与 $f\circ s'$ 也定义 $\phi$，这说明 $\phi$ 是 $D^b(\mathcal{M}_{fg}(U_\theta))$ 中的态射。这就证明了 $i$ 是全忠实的。又由 …… 可知它本质满。
\end{proof}

这个结果特别意味着：对任意 $V^\bullet\in D^b_{fg}(U_\theta)$，其局部化 $L\Delta_\lambda(V^\bullet)\in D^-(\D_\lambda)$ 是上同调凝聚的复形。为更精确地讨论这一点，先引入 $D^b(\D_\lambda)$ 的几个子范畴。由于 $\mathcal{M}_{coh}(\D_\lambda)$ 是 $\mathcal{M}(\D_\lambda)$ 的厚子范畴，可以定义范畴 $D^b_{coh}(\D_\lambda)$，即 $D^b(\D_\lambda)$ 中由上同调凝聚的有界复形组成的满子范畴，以及凝聚 $\D_\lambda$-模的导出范畴 $D^b(\mathcal{M}_{coh}(\D_\lambda))$。存在从 $D^b(\mathcal{M}_{coh}(\D_\lambda))$ 到 $D^b_{coh}(\D_\lambda)$ 的自然函子 $\Psi_\lambda$，它关于平移函子满足自然的交换性。

\begin{lemma}\label{lem:03.1.20}
自然函子 $\Psi_\lambda:D^b(\mathcal{M}_{coh}(\D_\lambda))\to D^b_{coh}(\D_\lambda)$ 是范畴等价。
\end{lemma}

\begin{proof}
利用平移函子可设 $\lambda$ 反支配且正则。由 \ref{cor:03.1.3} 与 \ref{thm:03.1.6}，此结果由 \ref{lem:03.1.19} 得出。
\end{proof}

因此，对正则的 $\lambda\in\h^*$，可以把函子 $L\Delta_\lambda$ 看作从 $D^b_{fg}(U_\theta)$ 到 $D^b_{coh}(\D_\lambda)$ 的函子。此外还有下面的结果。

\begin{theorem}\label{thm:03.1.21}
设 $\lambda\in\h^*$ 正则。则函子 $L\Delta_\lambda:D^b_{fg}(U_\theta)\to D^b_{coh}(\D_\lambda)$ 是范畴等价，其逆为 $R\Gamma$。
\end{theorem}

为证明这一陈述，只需说明：对任意 $V^\bullet\in D^b_{coh}(\D_\lambda)$ 有 $R\Gamma(V^\bullet)\in D^b_{fg}(U_\theta)$。这是 \ref{lem:03.1.17} 的下述推广的推论。先需要一个简单的引理。

\begin{lemma}\label{lem:03.1.22}
设 $V$ 是有限生成 $\g$-模，$F$ 是有限维 $\g$-模。则 $V\otimes_\C F$ 是有限生成 $\g$-模。
\end{lemma}

\begin{proof}
设 $U$ 是 $V$ 的有限维子空间，它作为 $\g$-模生成 $V$。我们断言 $U\otimes_\C F$ 作为 $\g$-模生成 $V\otimes_\C F$。设 $W$ 是 $V\otimes_\C F$ 中由 $U\otimes_\C F$ 生成的 $\g$-子模。定义 $V'$ 为 $V$ 中所有满足「对一切 $f\in F$ 有 $v\otimes f\in W$」的 $v$ 组成的子集。显然 $V'$ 是 $V$ 的线性子空间。又对任意 $\xi\in\g$、$v\in V'$ 与 $f\in F$ 有
$$\xi v\otimes f=\xi(v\otimes f)-v\otimes\xi f\in W,$$
故 $\xi v\in V'$，即 $V'$ 是 $V$ 的 $\g$-子模。另一方面，$V'$ 包含 $U$，故 $V'=V$。这又给出 $W=V\otimes_\C F$。
\end{proof}

\begin{lemma}\label{lem:03.1.23}
设 $\lambda\in\h^*$。则对任意 $\mathcal{V}\in\mathcal{M}_{coh}(\D_\lambda)$，$U_\theta$-模 $H^i(X,\mathcal{V})$（$i\in\Z_+$）都是有限生成的。
\end{lemma}

\begin{proof}
对 $i$ 作归纳，$0\le i\le\dim X$。可取支配权 $\mu$，使 $\lambda-\mu$ 反支配且正则。则 $\mathcal{V}$ 的平移 $\mathcal{V}(-\mu)$ 是凝聚 $\D_{\lambda-\mu}$-模；由 \ref{lem:03.1.17}，$\Gamma(X,\mathcal{V}(-\mu))$ 是有限生成 $\g$-模。设 $F$ 是以 $\mu$ 为最高权的不可约有限维 $\g$-模。此时对一切 $i$（$0\le i\le\dim X$）有
$$H^i(X,\mathcal{V}(-\mu)\otimes_{\calO_X}F)=H^i(X,\mathcal{V}(-\mu))\otimes_\C F,$$
故当 $i>0$ 时左端为零。另一方面，C.2 中研究的 $F$ 的滤过给出 $F_1=\calO(\mu)$ 到 $F$ 的单射。于是用 $\mathcal{V}(-\mu)$ 作张量积，得到 $U^\circ$-模的正合列
$$0\longrightarrow\mathcal{V}\longrightarrow\mathcal{V}(-\mu)\otimes_{\calO_X}F\longrightarrow K\longrightarrow 0.$$
对这个正合列作用 $\Gamma$，可见 $\Gamma(X,\mathcal{V})$ 是张量积 $\Gamma(X,\mathcal{V}(-\mu))\otimes_\C F$ 的 $\g$-子模，而后者由 \ref{lem:03.1.22} 有限生成。

\begin{annote}
底本此处作「by 20」，但 1.20 是导出范畴等价的陈述，与有限生成无关；此处所需的正是紧邻的 1.22（有限维 $\g$-模张量有限生成模仍有限生成），故按上下文改指 \ref{lem:03.1.22}。[来源：底本 p.42，上下文自证]
\end{annote}

这就证明了 $i=0$ 的情形。设结论对 $k-1$ 成立，$k\ge 1$。则由上同调的长正合列可知，$H^k(X,\mathcal{V})$ 是 $H^{k-1}(X,K)$ 的商。另一方面，由 $F$ 的滤过的定义可知，$K$ 带有自然的 $U^\circ$-模滤过，其相应的分次模 $\Gr K$ 等于 $\bigoplus_\nu\mathcal{V}(-\mu+\nu)$，其中和取遍 $F$ 的所有不等于 $\mu$ 的权 $\nu$。由归纳假设，$H^{k-1}(X,\mathcal{V}(-\mu+\nu))$ 都是有限生成 $\g$-模。对 $K$ 的滤过长度作归纳，即知 $H^{k-1}(X,K)$ 是有限生成 $\g$-模。
\end{proof}

最后，导出范畴的等价（\ref{thm:03.1.12}）与 Borel--Weil--Bott 定理（C.5.1）给出下述直接推论。

\begin{proposition}\label{prop:03.1.24}
设 $F$ 是以 $\lambda$ 为最低权的有限维不可约 $\g$-模。则对任意 $\mu=w(\lambda-\rho)$（$w\in W$）有：当 $p\ne-\ell(w)$ 时 $L^p\Delta_\mu(F)=0$，且
$$L^{-\ell(w)}\Delta_\mu(F)=\calO(\mu+\rho).$$
\end{proposition}
